\chapter{Interface générale}
\label{chap:interface}

\section{Vue d'ensemble}

La fenêtre principale d'AirfoilTools est structurée en quatre
zones (\figref{01_main_window.png}) :

\begin{enumerate}
    \item la \textbf{barre de menus} (\menu{Fichier}, \menu{Édition},
        \menu{Affichage}, \menu{Options}, \menu{Aide}) ;
    \item le \textbf{sélecteur d'onglets} (\menu{Profils},
        \menu{Paramétrage XFoil}, \menu{Résultats},
        \menu{Cp / Couche limite}) ;
    \item la \textbf{zone de travail} qui change selon l'onglet
        sélectionné ;
    \item la \textbf{barre de statut}, en bas de la fenêtre, où
        s'affichent les messages d'information et la progression
        des simulations.
\end{enumerate}

\section{Bulles d'aide}

Chaque contrôle (case à cocher, champ de saisie, bouton, action
de menu) dispose d'une \textbf{bulle d'aide} qui s'affiche au
survol prolongé du curseur souris. Ces bulles décrivent le rôle
du contrôle et les valeurs typiques.

\begin{astuce}
Survolez un contrôle dont la fonction n'est pas évidente :
la bulle d'aide apparaît au bout d'environ une seconde.
Pour les actions de menu, l'aide s'affiche en bas de
fenêtre dans la barre de statut au moment où l'on survole
l'entrée de menu.
\end{astuce}

\section{Menu \menu{Fichier}}

\begin{tabularx}{\linewidth}{l l X}
    \toprule
    \textbf{Entrée} & \textbf{Raccourci} & \textbf{Action} \\
    \midrule
    Ouvrir profil courant\dots
        & \touche{Ctrl}+\touche{O}
        & Charge un profil dans la position courante (bleu).
          Formats acceptés : \chemin{.dat} (Selig/Lednicer),
          \chemin{.bspl} (spline AirfoilTools), \chemin{.bez}
          (Bézier AirfoilTools), \chemin{.csv}.\\
    Ouvrir profil référence\dots
        & \touche{Ctrl}+\touche{Maj}+\touche{O}
        & Charge un profil dans la position de référence (rouge).\\
    Depuis la base UIUC / Profil courant\dots
        & ---
        & Ouvre un dialogue de recherche dans la base de données
          UIUC (Selig) en ligne et télécharge le profil sélectionné.
          Voir section~\ref{sec:uiuc}, page~\pageref{sec:uiuc}.\\
    Depuis la base UIUC / Profil référence\dots
        & ---
        & Idem pour le profil de référence.\\
    Profil NACA / Profil courant\dots
        & ---
        & Génère un profil NACA (4 ou 5 chiffres) comme profil courant.
          Demande les indices puis le nombre de points. Voir
          section~\ref{sec:naca}, page~\pageref{sec:naca}.\\
    Profil NACA / Profil référence\dots
        & ---
        & Idem pour le profil de référence.\\
    Sauvegarder profil\dots
        & \touche{Ctrl}+\touche{S}
        & Sauvegarde le profil courant. Un premier dialogue propose
          d'enregistrer le \emph{profil complet}, l'\emph{extrados}
          seul ou l'\emph{intrados} seul ; le format est ensuite
          sélectionné dans le filtre du dialogue de fichier.\\
    Sauvegarder profil avec volet\dots
        & ---
        & Sauvegarde le profil avec volet braqué (vert), de la même
          manière que le profil courant. Nécessite que le volet soit
          actif. Voir section~\ref{sec:flap}, page~\pageref{sec:flap}.\\
    Ouvrir projet\dots
        & \touche{Ctrl}+\touche{Maj}+\touche{P}
        & Ouvre un projet AirfoilTools (\chemin{.aftproj}) :
          restaure les profils courant et référence ainsi que
          l'image de calque et sa transformation. Voir
          chapitre~\ref{chap:formats}.\\
    Enregistrer projet\dots
        & \touche{Ctrl}+\touche{Alt}+\touche{S}
        & Enregistre l'ensemble de la session (profils + image de
          calque) dans un fichier \chemin{.aftproj}.\\
    Charger image de calque\dots
        & ---
        & Charge une image (jpg, png, tif, bmp) en arrière-plan pour
          décalquer un profil. Voir section~\ref{sec:calque},
          page~\pageref{sec:calque}.\\
    Retirer l'image de calque
        & ---
        & Supprime l'image de calque de la session courante.\\
    Quitter
        & \touche{Ctrl}+\touche{Q}
        & Ferme l'application.\\
    \bottomrule
\end{tabularx}

\section{Menu \menu{Édition}}

\begin{tabularx}{\linewidth}{l l X}
    \toprule
    \textbf{Entrée} & \textbf{Raccourci} & \textbf{Action} \\
    \midrule
    Convertir en Spline
        & \touche{Ctrl}+\touche{B}
        & Convertit le profil courant (mode points discrets)
          en mode spline multi-segment. Voir la section
          \emph{Conversion en spline}, page~\pageref{sec:convertspline}.\\
    Échantillonnage / Profil courant\dots
        & ---
        & Modifie le nombre de points d'échantillonnage du
          profil courant. Disponible uniquement en mode Spline.\\
    Échantillonnage / Profil référence\dots
        & ---
        & Idem pour le profil de référence.\\
    \bottomrule
\end{tabularx}

\section{Menu \menu{Affichage}}

\begin{tabularx}{\linewidth}{l l X}
    \toprule
    \textbf{Entrée} & \textbf{Raccourci} & \textbf{Action} \\
    \midrule
    Zoom adapté
        & \touche{Ctrl}+\touche{0}
        & Recadre la vue du canvas \menu{Profils} sur l'ensemble
          des profils visibles.\\
    Disposition / 1$\times$1, 1$\times$2, \dots
        & ---
        & Configure la grille de l'onglet \menu{Résultats}
          (lignes $\times$ colonnes).\\
    \bottomrule
\end{tabularx}

\section{Menu \menu{Options}}

\begin{tabularx}{\linewidth}{l l X}
    \toprule
    \textbf{Entrée} & \textbf{Raccourci} & \textbf{Action} \\
    \midrule
    Déviation / Verticale (épaisseur)
        & ---
        & Calcule la déviation entre profils comme un écart
          d'épaisseur mesuré \textbf{verticalement} (selon l'axe $z$),
          à abscisse constante. Mode par défaut.\\
    Déviation / Normale (perpendiculaire)
        & ---
        & Calcule la déviation \textbf{perpendiculairement} à la
          surface du profil courant (distance le long de la normale).
          Voir section~\ref{sec:deviation},
          page~\pageref{sec:deviation}.\\
    \bottomrule
\end{tabularx}

\begin{noteinfo}
Les deux modes de déviation sont \textbf{exclusifs} : choisir l'un
décoche automatiquement l'autre. Le changement est répercuté
immédiatement sur le tracé de l'onglet \menu{Profils}.
\end{noteinfo}

\section{Menu \menu{Aide}}

\begin{tabularx}{\linewidth}{l l X}
    \toprule
    \textbf{Entrée} & \textbf{Raccourci} & \textbf{Action} \\
    \midrule
    Manuel utilisateur
        & \touche{F1}
        & Ouvre le présent manuel (fichier PDF) avec le visualiseur
          par défaut du système. Le PDF est embarqué dans la
          distribution \chemin{.exe} et reste accessible même hors
          ligne.\\
    À propos\dots
        & ---
        & Affiche les informations de version et l'auteur.\\
    \bottomrule
\end{tabularx}

\begin{noteinfo}
En mode développement (lancement via \chemin{run\_gui.py}), le
manuel est lu depuis \chemin{docs/manuel/manuel.pdf} dans
l'arborescence du projet. En mode exécutable PyInstaller, il est
extrait depuis le bundle interne (dossier
\chemin{\_internal/docs/manuel.pdf} à côté de l'exe).
\end{noteinfo}

\section{Onglets}

Les quatre onglets sont indépendants : changer d'onglet ne perd
aucune donnée. La barre de statut peut afficher des messages
contextuels selon l'onglet actif. L'onglet \menu{Cp / Couche limite}
(chapitre~\ref{chap:cp}) offre une vue détaillée, façon \emph{XFoil},
d'un calcul unique.

\subsection*{Indicateur de résultats obsolètes}

L'onglet \menu{Résultats} possède un mécanisme particulier :
lorsqu'un profil est modifié après que la simulation a été lancée,
un \textbf{bandeau jaune} apparaît en haut de l'onglet pour signaler
que les résultats affichés ne correspondent plus aux profils
courants. Il est alors recommandé de relancer la simulation.

\begin{noteinfo}
Le bandeau d'obsolescence ne bloque pas la consultation des
résultats : il s'agit d'un simple indicateur visuel. La
suppression du bandeau se fait automatiquement au lancement
d'une nouvelle simulation.
\end{noteinfo}
